Currently, the large language models are the power behind of a large amount of everyday applications like chatbots, coding assistants, and media generation systems. Due to this big demand, most of cloud services providers are competing to make LLMs easily accessible in theirs platforms. However, running these models is far more expensive than running older search systems, because processing a single LLM request can cost up to ten times as much as a conventional keyword search~\cite{kumar2024latency}. Because of this, the method a server uses to decide which request to process next has a direct effect in the cloud costs resources utilization, and, of course, the speed of the output response.

In LLM serving systems, response speed is directly measured through latency, which represents the time a user waits from sending a prompt to receiving the complete generated output. When several users send queries simultaneously, incoming requests enter a server queue waiting to be processed. The component responsible for determining the order in which these queued tasks are executed is the scheduler, and currently, the most widely adopted baseline in production is First-Come-First-Served (FCFS)~\cite{yu2022orca}. Even though the FCFS handles requests in sequential order and avoids complex logic, it completely ignores task sizes, causing short requests to wait behind long ones. To improve the system latency, conventional scheduling algorithms theory often suggests the Shortest Job First (SJF) approach, because it prioritizes shorter tasks to speed up the execution time, but introduces the risk of starvation, leaving heavy requests waiting indefinitely. This is a classical scheduling problem to which Brinch Hansen proposed the best-known current solution: The response-ratio scheduling~\cite{hansen1973os}.

Because LLMs don't reveal the request sizes upfront, applying SJF requires predicting output token lengths before execution. Following this idea, recent works like~\cite{fu2024ltr} introduced Learning-to-Rank (LTR), a predictor-based scheduler that gets low latency by giving priority to requests with smaller predicted sizes. However, tests in~\cite{kumar2024latency} show that LTR causes severe starvation under high load and gets worse quickly when predictions fail or face request types it wasn't built for. FCFS finishes every request eventually, but with high latency. LTR works the opposite way: it cuts average wait times, but fairness suffers. Combining predictive scheduling with a mechanism that keeps starvation close to zero feels like a much more balanced option to explore.

The main goal of this project was to translate this idea of a balanced alternative into a practical scheduler, which led us to create a formula that scores each request using both its predicted size and its waiting time. However, running these tests on real GPU servers would not be possible because of high costs and limited hardware available, so using a discrete-event simulator was the solution found for this research. Recent tools like Vidur~\cite{agrawal2024vidur} show that a custom simulator can model LLM queues well without needing huge computational resources. Based on this, we built a discrete-event simulator in Python to compare our scoring formula against FCFS~\cite{yu2022orca} and LTR~\cite{kumar2024latency, fu2024ltr}, testing if this approach can keep low latency without causing request starvation. 

Our simulator runs a single-server queue using synthetic workloads as well as real prompt lengths from the LMSYS dataset~\cite{zheng2024lmsys}. This allowed us to see how the scheduler works under different traffic levels and with fake prediction errors that act like real-world predictor failures. This work makes the following contributions:

\begin{itemize}
    \item A priority formula that balances the prediction uncertainty with the waiting time that a request already spent in the queue.
    \item A parameter-selection method validated across all the tested conditions, preventing setups that works only in one specific case.
    \item An evaluation approach that considers also the requests that the scheduler fails to complete, and not only the ones that it finishes.
    \item A harder noise model created for testing the scheduler behavior under realistic failures beyond simple random errors.
\end{itemize}